\chapter{Apresentação da 1ª edição}

O Conselho Regional de Psicologia de São Paulo --- 6ª Região (CRP~SP),
entidade dotada de personalidade jurídica de direito público, com
autonomia administrativa e financeira, nos termos da Lei Federal
nº~5.766, de 20 de dezembro de 1971, tem como finalidade fiscalizar o
exercício da profissão de psicóloga/o, competindo-lhe orientar,
disciplinar e zelar pela fiel observância dos princípios
ético-profissionais e contribuir para o desenvolvimento da Psicologia
enquanto ciência e profissão.

Este CRP~SP tem sede na cidade de São Paulo e jurisdição no estado de
São Paulo, possuindo subsedes nas regiões de Alto Tietê (Mogi das
Cruzes), Assis (Assis), Baixada Santista e Vale do Ribeira (Santos),
Bauru (Bauru), Campinas (Campinas), Grande ABC (Santo André),
Metropolitana (São Paulo), Ribeirão Preto (Ribeirão Preto), São José do
Rio Preto (São José do Rio Preto), Sorocaba (Sorocaba) e Vale do Paraíba
e Litoral Norte (Taubaté).

O CRP~SP tem como atribuições principais: adotar as medidas e
procedimentos necessários à permanente orientação, disciplina e
fiscalização do exercício da profissão de psicóloga/o no estado de São
Paulo; adotar medidas e procedimentos para preservação do livre
exercício da profissão de psicóloga/o, bem como do respeito às suas
prerrogativas e direitos profissionais; executar os serviços
concernentes ao registro profissional das/os psicólogas/os, realizando
as inscrições, reativações, transferências e cancelamentos de registros,
expedindo às/aos inscritas/os a carteira de identidade profissional
(CIP); funcionar como tribunal regional de ética profissional; elaborar
e encaminhar ao Conselho Federal de Psicologia (CFP) a proposta
orçamentária anual, além do relatório geral de suas atividades;
providenciar a instalação da Assembleia Geral das/os psicólogas/os
inscritas/os na região; arrecadar anuidades, taxas e demais rendimentos
que lhe competem, promovendo o repasse da arrecadação ao Conselho
Federal; criar os serviços necessários ao bom desempenho de suas
atividades e autorizar a compra de material para suas instalações;
organizar o quadro de pessoal; zelar pela dignidade, independência,
prerrogativas e valorização da Psicologia; e zelar pela gestão
responsável, cumprindo a legislação a partir dos princípios da
administração pública: legalidade, impessoalidade, moralidade,
publicidade e eficiência.

Assim, para uma adequada utilização dos recursos e dispositivos
públicos, foi aprovado, no Planejamento Estratégico 2023--2025,
resultado específico para garantir a implantação adequada dos sistemas e
instrumentos oficiais do CRP~SP.

O Resultado 1.2 do Planejamento Estratégico supracitado se refere a
``ter implementado estrutura de gestão democrática com processos de
trabalho planejados e institucionalizados, de forma transversal,
acessível, integrada, transparente e com produção de dados.''

Das ações previstas para 2023, foi atingido o Resultado 1.2: ``Ter
implantado e integrado os sistemas informatizados (SEI, BRC, Benner e
produtos Zimbra/Implanta), garantindo o cumprimento das normativas
(leis, resoluções, Corep), dos fluxos e processos de trabalho com
formação permanente, transparência, produção, análise e segurança de
dados.''

E, por fim, fazendo cumprir a Macroação nº~8 --- ``Ter elaborado
política que institucionaliza e da segurança ao uso dos novos sistemas''
---, é elaborada esta Política de Segurança em Tecnologia da Informação
e Comunicação (PSTIC), que especifica os controles internos aplicáveis à
segurança, ao sigilo da informação e ao uso das ferramentas do Conselho
Regional de Psicologia da 6ª Região --- CRP~SP, com o objetivo de
realizar suas operações gerais, ainda que em situações adversas.

Estão sujeitas/os ao disposto no presente documento todas/os as/os
conselheiros, colaboradoras/es, trabalhadoras/es, funcionárias/os
terceirizadas/os, estagiárias/os, prestadoras/es de serviços e demais
convidadas/os, no que a cada um for aplicável.

A Política Segurança da TIC é o documento que normatiza, orienta e
estabelece as diretrizes para a proteção dos sistemas, documentos,
operações e comunicação institucional quanto a processos de trabalho e
uso da rede e, principalmente com a prevenção de responsabilidade legal
para todos as/os usuárias/os.

Além disso estabelecimento de fluxos para comunicação e guarda de
documentos e deve, portanto, ser cumprida e aplicada em todas as áreas
da autarquia.

A presente política está baseada nas leis vigentes brasileiras que
tangem sobre isso, bem como nas recomendações propostas pelas ABNT NBR
ISO/IEC 27002/2022, reconhecida mundialmente como um código de prática
para a gestão da segurança da informação.

Também é importante destacar que esta política visa apoiar inicialmente
a Lei Geral de Proteção dos Dados --- LGPD (Lei nº~13.709/2018) no qual
toda a informação produzida ou recebida como resultado da atividade
institucional, pertence ao CRP~SP.

Por fim, garantimos aqui que os registros institucionais sejam
disponibilizados no Portal da Transparência do CRP~SP preconizados pela
Lei de Acesso à Informação (Lei nº~12.527/2011) tomadas as devidas
proteções ao sigilo daqueles documentos internos que ficarão guardados
pela Autarquia.

Talita Fabiano de Carvalho

Conselheira presidenta do CRP-06